\chapter{Asynchronous Messaging with Pub/Sub}
\index{Pub/Sub}\index{topics}\index{subscriptions}\index{dead-letter topic}\index{snapshot and seek}\index{message schema}

\begin{objectives}
\begin{itemize}
    \item Model Pub/Sub topics, subscriptions (push and pull), and acknowledgements correctly.
    \item Enforce data contracts at the edge with Avro or Protocol Buffer topic schemas.
    \item Publish resiliently with retry and exponential backoff, and deduplicate at-least-once delivery downstream.
    \item Configure dead-letter topics, retry policies, and snapshot/seek replay.
    \item Compare Pub/Sub Standard and Pub/Sub Lite on cost and operational overhead.
    \item Architect a global telemetry ingestion point for five million mobile games with exactly-once processing semantics.
\end{itemize}
\end{objectives}

\begin{crosscloud}
Managed messaging services share one mental model---producers publish to named channels, consumers subscribe, and the broker owns durability and delivery---with different knobs for ordering, replay, and cost.

\vspace{2mm}
{\footnotesize
\begin{tabular}{@{}p{2.8cm}p{3.0cm}p{3.0cm}p{3.0cm}@{}} \\
\toprule
\textbf{Capability} & \textbf{GCP Pub/Sub} & \textbf{AWS} & \textbf{Azure} \\
\midrule
Managed pub/sub & Pub/Sub (global) & SNS + SQS / EventBridge & Service Bus / Event Grid \\
Kafka-compatible & Managed Service for Apache Kafka & MSK & Event Hubs (Kafka surface) \\
Dead-lettering & Dead-letter topics & SQS DLQs / EventBridge DLQ & Service Bus dead-letter queues \\
Replay & Snapshots + seek & (archive + replay in EventBridge) & (replay via offsets on Event Hubs) \\
Schema enforcement & Avro / Proto topic schemas & EventBridge schemas & Schema registry (Event Hubs) \\
Cost-profiled tier & Pub/Sub Lite (zonal, capacity) & Kinesis on-demand vs provisioned & Event Hubs throughput units \\
\bottomrule
\end{tabular}}

\vspace{1mm}
Self-managed \textbf{Apache Kafka} remains the industry reference point for partition-based, replayable logs \cite{apachekafkadocs}. Pub/Sub differs deliberately: global by default, partition-free for the user, and priced per GiB delivered rather than per provisioned broker.
\end{crosscloud}

\begin{keyterms}
\textbf{Topic}: a named channel to which publishers send messages.
\textbf{Subscription}: a named cursor over a topic delivering messages to consumers by push or pull.
\textbf{Acknowledgement (ack)}: the consumer's confirmation that a message was processed; unacked messages are redelivered.
\textbf{Dead-letter topic}: a topic receiving messages that exceeded a subscription's maximum delivery attempts.
\textbf{Snapshot/seek}: Pub/Sub's replay mechanism, rewinding a subscription to a retained point in time.
\textbf{Exactly-once delivery}: a subscription mode where Pub/Sub itself suppresses redelivery of already-acked messages within a region.
\end{keyterms}

\section{Week 13 Roadmap}
Week~13 opens Part~IV with the front door of every streaming platform. Monday covers topics, subscriptions, acknowledgements, and schema enforcement at the edge. Tuesday covers failure engineering: dead-letter topics, retry policies, snapshots, and seek. Wednesday compares Pub/Sub Lite with Standard. Thursday publishes mock financial transactions with a dead-letter safety net. Friday architects global game telemetry with exactly-once semantics.

\begin{definitionbox}
\textbf{Pub/Sub} is Google Cloud's global, serverless messaging service. Publishers send messages to topics; subscriptions deliver them to consumers with at-least-once semantics by default, message retention, flow control, and optional exactly-once delivery and ordering keys \cite{googlecloudpubsubdocs}.
\end{definitionbox}

\section{Monday: Topics, Subscriptions, Acks, and Schema Enforcement}
\index{Pub/Sub!topics}\index{Pub/Sub!subscriptions}\index{acknowledgement}\index{topic schema}

\subsection{The Delivery Model}
A topic is global and durable; messages are retained for a configurable window regardless of consumption. Each subscription is an independent cursor: ten subscriptions on one topic each receive every message---fan-out is free, which is why a single ingestion topic commonly feeds a serving pipeline, an analytics sink, and an audit archive simultaneously. Pull subscriptions have consumers request batches; push subscriptions deliver to an HTTPS endpoint with OIDC authentication---push fits serverless consumers (Cloud Run, Cloud Functions) that should not run polling loops.

\subsection{Acknowledgement Semantics}
Delivery is at-least-once: Pub/Sub redelivers any message not acked within the subscription's acknowledgement deadline. Consumers must therefore be idempotent or deduplicate. Ack deadlines can be extended per message (\texttt{modifyAckDeadline}) while processing continues; abandoning extensions is how a slow consumer generates its own duplicate storm.

\subsection{Schema Validation at the Edge}
Attaching an Avro or Protocol Buffer schema to a topic makes the broker reject non-conforming publishes: the data contract enforced at the edge, before a malformed message can reach any of the ten subscriptions. Schema evolution follows compatibility rules (additive fields); versioning lives on the schema, and breaking changes mean a new schema revision with a deployment plan.

\begin{pitfall}
Schema enforcement protects message \emph{shape}, not message \emph{semantics}. A schema-valid message can still carry a negative price or a timestamp from 1970. Edge validation is the first contract layer, not the only one---validate semantics in the pipeline (Chapters 14--15) and at load (Chapter~11).
\end{pitfall}

\begin{tipbox}
Publish with attributes (\texttt{origin}, \texttt{schema\_version}, \texttt{event\_type}) even when the payload is self-describing. Attributes are filterable on subscriptions and visible in dead-letter triage without deserializing the body.
\end{tipbox}

\section{Tuesday: Dead Letters, Retries, Snapshots, and Seek}
\index{dead-letter topic!configuration}\index{retry policy}\index{snapshot}\index{seek}

\subsection{Dead-Letter Topics}
Configure a dead-letter topic \emph{on the subscription} with a maximum delivery attempt count: after N failed deliveries the message forwards to the DLQ with its attempt metadata attached. The DLQ is a first-class topic---subscribe an alerting pipeline to it, sample its messages daily, and own the triage runbook. A DLQ nobody reads is a data-loss mechanism with good branding.

\subsection{Retry Policies and Backoff}
Subscription retry policies add exponential backoff between redeliveries (minimum and maximum delay), smoothing thundering-herb redelivery against a struggling consumer. Publisher-side, the client libraries retry transient failures with their own backoff; both layers compose, so reason about worst-case latency of a message that retries on both sides.

\subsection{Snapshots and Seek}
\textbf{Snapshots} capture a subscription's ack state; \textbf{seek} rewinds a subscription to a snapshot or a timestamp within the topic retention window. Seek is the streaming answer to ``the consumer had a bug for six hours'': fix the consumer, seek back to before the bad deploy, replay. It is also the migration tool for consumer changes that must reprocess history.

\begin{lstlisting}[language=bash,caption={DLQ configuration and a seek-based replay from PowerShell.}]
$env:PROJECT_ID = "my-gcp-pubsub-lab"

# DLQ is configured on the subscription, not the publisher:
gcloud pubsub subscriptions create card-sub `
  --topic=card-transactions `
  --dead-letter-topic=card-dlq `
  --max-delivery-attempts=5 `
  --min-retry-delay=10s --max-retry-delay=600s

# Replay: seek the subscription back two hours after fixing a consumer bug.
gcloud pubsub subscriptions seek card-sub --time=$((Get-Date).ToUniversalTime().AddHours(-2).ToString('yyyy-MM-ddTHH:mm:ssZ'))
\end{lstlisting}

\begin{notebox}
Seek replays within retention. A topic retaining one day cannot replay last week's incident. Size retention to your recovery-time expectations, not to the default.
\end{notebox}

\section{Wednesday: Pub/Sub Lite versus Pub/Sub Standard}
\index{Pub/Sub Lite}\index{cost vs operational overhead}

Pub/Sub Lite trades Standard's global, fully-managed elasticity for zonal, capacity-provisioned, partition-based topics at substantially lower per-GiB cost. Lite requires you to reason about partitions, throughput reservations, and zonal availability; Standard reasons about none of them. The decision rule: Standard by default; Lite when message volume is large and steady, latency tolerance includes zonal failure modes, and the team can operate capacity planning. Many platforms use both: Standard for control-plane and low-volume topics, Lite for the firehose.

\begin{table}[H]
\centering
\small
\begin{tabular}{@{}p{3.6cm}p{4.4cm}p{4.4cm}@{}} \\
\toprule
\textbf{Dimension} & \textbf{Pub/Sub Standard} & \textbf{Pub/Sub Lite} \\
\midrule
Scope & Global, multi-zone durability & Zonal (Lite) or regional \\
Pricing & Per GiB delivered + retention & Provisioned capacity, lower per-GiB \\
Partitions & Abstracted away & Explicit; ordering per partition \\
Capacity & Elastic & Reserved throughput \\
Best for & Default; variable workloads & High steady volume, cost-sensitive \\
\bottomrule
\end{tabular}
\caption{Pub/Sub Standard versus Lite \cite{googlecloudpubsubdocs}.}
\label{tab:ch13-lite-vs-standard}
\end{table}

\section{Thursday Hands-On Project: Financial Transactions with a Dead-Letter Net}
\index{hands-on project}\index{publisher retry}\index{exponential backoff}
This lab creates a topic and pull subscription, publishes one thousand mock card transactions from Python with retry and backoff, and routes poison messages to a dead-letter topic.

\subsection{Provision}
\begin{lstlisting}[language=bash,caption={Creating topic, DLQ, and subscription with delivery guardrails.}]
$env:PROJECT_ID = "my-gcp-pubsub-lab"
gcloud config set project $env:PROJECT_ID
gcloud services enable pubsub.googleapis.com

gcloud pubsub topics create card-transactions
gcloud pubsub topics create card-dlq

gcloud pubsub subscriptions create card-sub `
  --topic=card-transactions `
  --dead-letter-topic=card-dlq `
  --max-delivery-attempts=5 `
  --ack-deadline=60
\end{lstlisting}

\subsection{Publish with Backoff}
\begin{lstlisting}[language=Python,caption={Publishing mock card transactions with retry and exponential backoff.}]
import json, random, time
from google.cloud import pubsub_v1

publisher = pubsub_v1.PublisherClient(
    publisher_options=pubsub_v1.types.PublisherOptions(
        retry=pubsub_v1.types.Retry(initial=0.1, maximum=10.0, multiplier=2.0)))
topic_path = publisher.topic_path("my-project", "card-transactions")

def publish(txn):
    future = publisher.publish(topic_path, json.dumps(txn).encode("utf-8"),
                               origin="pos-simulator")
    future.add_done_callback(lambda f: print("acked:", f.result()))

for i in range(1000):
    publish({"card_id": f"card-{random.randint(1, 500)}",
             "amount": round(random.uniform(1, 500), 2),
             "city": random.choice(["NYC", "LON", "TYO"]), "ts": time.time()})
\end{lstlisting}

\subsection{Consume, Fail, and Inspect the DLQ}
Write a pull consumer that acks valid messages and deliberately nacks (or lets the ack deadline expire on) messages with \texttt{amount > 490} to simulate a consumer bug. After five attempts those messages should land in \texttt{card-dlq}; pull from a DLQ subscription to verify, and inspect the \texttt{delivery\_attempt} attribute.

\subsection*{Deliverables}
\begin{itemize}
    \item Publisher and consumer scripts with a README describing retry and DLQ behavior.
    \item Evidence of poison messages arriving in \texttt{card-dlq} with delivery-attempt metadata.
    \item A screenshot of subscription metrics: oldest unacked message age and ack rate during the run.
\end{itemize}

\subsection*{Acceptance Criteria}
\begin{itemize}
    \item The publisher retries transient failures with exponential backoff.
    \item The DLQ receives exactly the messages that exhaust delivery attempts---no more, no fewer.
    \item The consumer is idempotent: redelivered messages do not double-apply effects.
\end{itemize}

\subsection*{Stretch Goals}
\begin{itemize}
    \item Attach an Avro schema to the topic and show a non-conforming publish rejected.
    \item Enable exactly-once delivery on the subscription and describe what changes.
    \item Seek the subscription back fifteen minutes and observe redelivery of acked messages.
\end{itemize}

\section{Friday Use Case and Architecture: Telemetry for Five Million Mobile Games}
\index{mobile telemetry}\index{exactly-once processing}\index{global ingestion}

\subsection{Scenario}
A games publisher ingests session telemetry from five million concurrent mobile clients worldwide: events per second in the hundreds of thousands at peak, highly variable per region, with fraud, analytics, and live-ops consumers. The requirement reads ``exactly-once processing''---which must be decomposed before it can be built.

\subsection{Decomposing Exactly-Once}
No system delivers exactly-once \emph{delivery} across arbitrary failures; what platforms build is exactly-once \emph{effect}. The chain: publishers retry with backoff (duplicates possible); Pub/Sub deduplicates within its window and offers exactly-once delivery mode within a region; the Dataflow consumer (Chapters 14--15) uses Streaming Engine's exactly-once shuffle; the sink is BigQuery's Storage Write API in exactly-once mode or an idempotent key-based write. Each hop removes a duplicate class; the composition yields exactly-once effect end to end \cite{akidau2015dataflow}.

\begin{figure}[H]
    \centering
    \resizebox{0.95\textwidth}{!}{%
    \begin{tikzpicture}[
        src/.style={rectangle, rounded corners, draw=codegray, thick, fill=white, text width=2.7cm, minimum height=1.0cm, align=center, font=\small\bfseries},
        svc/.style={rectangle, rounded corners, draw=primary, thick, fill=primary!6, text width=2.9cm, minimum height=1.0cm, align=center, font=\small\bfseries},
        dst/.style={rectangle, rounded corners, draw=codegreen!60!black, thick, fill=green!7, text width=2.9cm, minimum height=1.0cm, align=center, font=\small\bfseries},
        arrow/.style={-{Stealth[length=2.5mm]}, draw=primary, thick},
        dasharrow/.style={-{Stealth[length=2.5mm]}, draw=codegray, thick, dashed}
    ]
        \node[src] (games) at (0,0) {5M game\\clients};
        \node[svc] (ps) at (3.8,0) {Pub/Sub\\global topics\\schema-enforced};
        \node[svc] (df) at (7.6,0) {Dataflow\\exactly-once\\shuffle};
        \node[dst] (bq) at (11.4,1.2) {BigQuery\\Storage Write\\(exactly-once)};
        \node[dst] (bt) at (11.4,-1.2) {Bigtable\\live-ops state};

        \draw[arrow] (games) -- (ps);
        \draw[arrow] (ps) -- (df);
        \draw[arrow] (df) -- (bq);
        \draw[arrow] (df) -- (bt);
    \end{tikzpicture}%
    }
    \caption{Global game telemetry: edge schema enforcement, exactly-once shuffle, idempotent sinks.}
    \label{fig:ch13-game-telemetry}
\end{figure}

\subsection{Operational Reality}
The interesting failures are not message loss---retention and redelivery cover that---but backlog: alert on oldest unacked message age and backlog bytes per subscription, because a consumer that falls behind at global scale catches up slowly. Per-region traffic shaping is free via subscription flow control; per-tenant abuse is handled by publisher-side quotas and schema/attribute filtering rather than by the broker.

\begin{pitfall}
Exactly-once delivery mode is regional and changes redelivery behavior; combining it casually with ordering keys and high fan-out can surprise throughput. Validate the mode under a load test that includes consumer crashes---the scenario it exists for.
\end{pitfall}

\section{Interview Q\&A}
\subsection{Concepts and Trade-offs}
\begin{enumerate}
    \item \textbf{Q: What happens to a message that exceeds max delivery attempts?}\\
    A: It forwards to the subscription's dead-letter topic with attempt metadata, where a triage process owns it---it is not silently deleted.

    \item \textbf{Q: How does a topic schema enforce a data contract at publish time?}\\
    A: The broker validates every published message against the Avro/Proto schema and rejects non-conforming publishes before any subscriber sees them.

    \item \textbf{Q: Why is Pub/Sub publishing at-least-once, and how do you deduplicate?}\\
    A: Acknowledgements can be lost after successful processing, forcing redelivery; deduplicate with message IDs, idempotent consumers, Dataflow exactly-once mode, or exactly-once sinks.

    \item \textbf{Q: When does Pub/Sub Lite beat Standard on cost?}\\
    A: At high, steady throughput where provisioned zonal capacity is cheaper per GiB than Standard's elastic per-GiB pricing and the team can operate partitions.

    \item \textbf{Q: What does the oldest-unacked-message-age metric warn you about?}\\
    A: Consumer lag or stuck messages: a growing value means the pipeline is falling behind or specific messages are never being acked (poison messages heading for the DLQ).
\end{enumerate}

\section{Further Reading}
Read the Pub/Sub documentation for subscriptions, schemas, exactly-once delivery, and Lite \cite{googlecloudpubsubdocs}. The Kafka documentation provides the partition-log comparison point \cite{apachekafkadocs}. The Dataflow Model paper is the theoretical foundation for the exactly-once chain assembled in the Friday use case \cite{akidau2015dataflow}, and Cloud Monitoring documentation covers the backlog alerting \cite{googlecloudmonitoringdocs}.

\section*{Key Takeaways}
\begin{itemize}
    \item Fan-out is free: model one ingestion topic with many subscriptions rather than many producer paths.
    \item Edge schemas reject malformed messages before they reach any consumer---shape first, semantics next.
    \item At-least-once delivery makes idempotent consumers and deduplication a requirement, not an optimization.
    \item DLQs, retry policies, and seek are the failure-engineering kit; a DLQ without triage is data loss.
    \item Pub/Sub Lite trades operational simplicity for cost at high steady volume.
    \item Exactly-once is a property of the whole chain---publisher, broker, shuffle, sink---assembled deliberately.
\end{itemize}

\section{Exercises}
\begin{enumerate}
    \item \textbf{(Conceptual)} Walk a single message through publish, delivery, ack-loss, redelivery, and deduplication, naming the duplicate-suppression mechanism at each hop.
    \item \textbf{(Conceptual)} A product manager demands ``no duplicate analytics events, ever.'' Explain what is and is not achievable, and the design that satisfies the intent.
    \item \textbf{(Hands-on)} Reproduce the Thursday lab, then raise \texttt{max-delivery-attempts} and observe the latency before DLQ arrival.
    \item \textbf{(Hands-on)} Attach a Proto schema, publish a non-conforming message, and capture the rejection error.
    \item \textbf{(Design)} Choose Lite versus Standard for three workloads: 200 msg/s control events, 2 GB/s steady telemetry, and 50 GB/s Black-Friday bursts.
    \item \textbf{(Architecture)} Add a replay drill to the telemetry design: consumer bug deployed for six hours, detected, fixed, and replayed via seek---write the runbook.
\end{enumerate}
